\section{The \texttt{Union\_mesh} McStas Component}
Component for including 3D mesh in Union geometry

\subsection*{Identification}
\begin{itemize}
  \item \textbf{Site:} 
  \item \textbf{Author:} Martin Olsen, and expanded upon by Daniel Lomholt Christensen
  \item \textbf{Origin:} University of Copenhagen / ACTNXT project
  \item \textbf{Date:} 17.09.18
\end{itemize}

\subsection*{Description}
\begin{lstlisting}
Part of the Union components, a set of components that work together and thus
separates geometry and physics within McStas.
The use of this component requires other components to be used.

1) One specifies a number of processes using process components
2) These are gathered into material definitions using Union_make_material
3) Geometries are placed using Union_box/cylinder/sphere/mesh, assigned a material
4) A Union_master component placed after all the above

Only in step 4 will any simulation happen, and per default all geometries
defined before this master, but after the previous will be simulated here.

There is a dedicated manual available for the Union_components

The mesh component loads a 3D off or stl files as the geometry.
The mesh geometry that is loaded must be a watertight mesh.
In order to check this, the component implements a simple Euler Pointcare value,
To check if the given geometry is watertight. This check is sometimes too rigid,
so a user can set the skip_convex_check parameter in order to not have this
check performed.


If you load an off file, we currently only support rank 3 polygons, i.e
triangular meshes.

It is allowed to overlap components, but it is not allowed to have two
parallel planes that coincides. This will crash the code on run time.
\end{lstlisting}

\subsection*{Input parameters}
Parameters in \textbf{boldface} are required; the others are optional.

\begin{longtable}{p{0.22\textwidth}p{0.12\textwidth}p{0.46\textwidth}p{0.14\textwidth}}
\toprule
\textbf{Name} & \textbf{Unit} & \textbf{Description} & \textbf{Default} \\
\midrule
\endhead
filename & str & Name of file that contains the 3D geometry. Can be either .OFF, .off, .STL, or .stl & 0 \\
material\_string & string & material name of this volume, defined using Union\_make\_material & 0 \\
\textbf{priority} & 1 & priotiry of the volume (can not be the same as another volume) A high priority is on top of low. &  \\
visualize & 1 & set to 0 if you wish to hide this geometry in mcdisplay & 1 \\
all\_surfaces &  &  & 0 \\
cut\_surface &  &  & 0 \\
target\_index & 1 & Focuses on component a component this many steps further in the component sequence & 0 \\
target\_x & m &  & 0 \\
target\_y & m & Position of target to focus at & 0 \\
target\_z & m &  & 0 \\
focus\_aw & deg & horiz. angular dimension of a rectangular area & 0 \\
focus\_ah & deg & vert. angular dimension of a rectangular area & 0 \\
focus\_xw & m & horiz. dimension of a rectangular area & 0 \\
focus\_xh & m & vert. dimension of a rectangular area & 0 \\
focus\_r & m & focusing on circle with this radius & 0 \\
p\_interact & 1 & probability to interact with this geometry [0-1] & 0 \\
mask\_string & string & Comma seperated list of geometry names which this geometry should mask & 0 \\
mask\_setting & string & "All" or "Any", should the masked volume be simulated when the ray is in just one mask, or all. & 0 \\
number\_of\_activations & 1 & Number of subsequent Union\_master components that will simulate this geometry & 1 \\
skip\_convex\_check & 1 & when set to 1, skips the check for whether input .off file is convex. & 0 \\
coordinate\_scale & 1e-3 & Multiplier that the vertex positions are multiplied by. Set to 1 for input coordinates in meters. Set to 1e-3 for input coordinates in mm. & 1e-3 \\
init & string & name of Union\_init component (typically "init", default) & "init" \\
\bottomrule
\end{longtable}

\subsection*{Links}
\begin{itemize}
  \item \href{run:/home/willend/willend-McCode/mcstas-comps/union/Union_mesh.comp}{Source code} for \texttt{Union\_mesh.comp}.
\end{itemize}
\IfFileExists{Union_mesh_static.tex}{\input{Union_mesh_static.tex}}{}